% !TeX root = ../despliegue.tex
%=======================================================================
%  config/plantilla.tex — mecanismos propios de este manual
%=======================================================================

%------------------------- Bloques de comandos --------------------------
% Este documento es, en buena parte, comandos que alguien va a copiar y
% pegar en una terminal. Se componen con fancyvrb dentro de una tcolorbox:
% fancyvrb —a diferencia de listings— compone sin problema los acentos de
% los comentarios bajo inputenc + T1, y la caja aporta el fondo gris.
%
% \VerbatimEnvironment es obligatorio: sin él, envolver un entorno
% verbatim dentro de \newenvironment no funciona.
\definecolor{codeback}{HTML}{F6F8FA}
\definecolor{codeframe}{HTML}{D0D7DE}

\newenvironment{comando}
  {\VerbatimEnvironment
   \begin{tcolorbox}[enhanced,breakable,
       colback=codeback,colframe=codeframe,boxrule=0.6pt,arc=2pt,
       left=8pt,right=8pt,top=6pt,bottom=6pt]
   \begin{Verbatim}[fontsize=\small,xleftmargin=0pt]}
  {\end{Verbatim}\end{tcolorbox}}

%--------------------------- Cajas de aviso -----------------------------
\definecolor{notaback}{HTML}{EFF6FF}
\definecolor{notaframe}{HTML}{1168BD}
\definecolor{avisoback}{HTML}{FEF3C7}
\definecolor{avisoframe}{HTML}{B45309}

% \nota{...} — información útil o aclaración, sin urgencia.
\newcommand{\nota}[1]{%
  \begin{tcolorbox}[enhanced,breakable,
      colback=notaback,colframe=notaframe,boxrule=0.6pt,arc=2pt,
      left=8pt,right=8pt,top=5pt,bottom=5pt,
      title={Nota},fonttitle=\bfseries\small,
      coltitle=notaframe,colbacktitle=notaback,
      fontupper=\small]
    #1
  \end{tcolorbox}}

% \aviso{...} — advertencia sobre algo irreversible o que requiere cuidado.
\newcommand{\aviso}[1]{%
  \begin{tcolorbox}[enhanced,breakable,
      colback=avisoback,colframe=avisoframe,boxrule=0.6pt,arc=2pt,
      left=8pt,right=8pt,top=5pt,bottom=5pt,
      title={Atención},fonttitle=\bfseries\small,
      coltitle=avisoframe,colbacktitle=avisoback,
      fontupper=\small]
    #1
  \end{tcolorbox}}

%--------------------------- Marcas en el texto -------------------------
% \cmd{docker compose up} — comando o nombre técnico dentro de un párrafo.
% \detokenize deja pasar el guion bajo de POSTGRES_PASSWORD y compañía sin
% tener que escaparlo en cada aparición.
\newcommand{\cmd}[1]{\texttt{\small\detokenize{#1}}}
% \var{GATEWAY_PORT} — variable de entorno.
\newcommand{\var}[1]{\texttt{\small\detokenize{#1}}}
% \servicio{ms-reservas} — nombre de un contenedor o servicio del compose.
\newcommand{\servicio}[1]{\texttt{\small\detokenize{#1}}}
